% Manuscript skeleton — Digital lifestyles and consumption-based carbon emissions
% Compile with pdflatex (Overleaf-compatible). Venue undecided; plain article for now.
\documentclass[11pt,a4paper]{article}

\usepackage[utf8]{inputenc}
\usepackage[T1]{fontenc}
\usepackage[english]{babel}
\usepackage[margin=2.5cm]{geometry}
\usepackage{setspace}
\usepackage{graphicx}
\usepackage{booktabs}
\usepackage{amsmath}
\usepackage{natbib}
\usepackage{xcolor}
\usepackage[hidelinks]{hyperref}

\graphicspath{{../output/}}

\newcommand{\todo}[1]{\textcolor{red}{[TODO: #1]}}

\title{Digital lifestyles and consumption-based carbon emissions:\\
individual-level evidence from linked bank-transaction, survey, and register data}
\author{David Andersson \and Mathias Lehner% \and G\"oran Finnveden \and Anna Furberg
}
\date{Draft --- \today}

\begin{document}

\maketitle

\begin{abstract}
\todo{Condense the extended abstract (latex/extended\_abstract.tex) to
$\sim$150--250 words once the venue is decided.}
\end{abstract}

\onehalfspacing

\section{Introduction}

\todo{Framing: optimists (dematerialization, decarbonization, demobilization) vs.\
skeptics (rebound, e-commerce-induced demand). Gap: no study observes digital
behavior and measured consumption-based emissions in the same individuals.}

\section{Data}

\subsection{Sample}
\todo{$\sim$6{,}000 Swedish adults (Konsumtionskollen sample); surveys 2024;
bank-account linkage retrieving transactions back to 2019.}

\subsection{Consumption-based carbon footprints}
\todo{Bank transactions categorized to COICOP and converted to
CO\textsubscript{2}e; person-level annual aggregation; P99 winsorization.}

\subsection{Digital lifestyle intensity}
\todo{Pre-registered index: device-use frequencies (baseline q11\_1--q11\_5+)
and digital-activity frequencies (q11b\_1--q11b\_8), anchored cardinally by
the device-assisted screen-time report (endline E4: iOS Sk\"armtid / Android
Digitalt v\"alm\aa ende, average daily h:min previous week). Report reliability
(Cronbach's $\alpha$). Note E4 opt-outs and missingness.}

\subsection{Register variables}
\todo{Income, education, household composition, DeSO-level urbanity.}

\section{Empirical strategy}

\subsection{Net gradient}
\todo{Total CO\textsubscript{2}e $\sim$ screen-time hours + sociodemographic
controls; HC3 robust SEs.}

\subsection{Decomposition}
\todo{Category-level gradients: transport (demobilization), goods and
e-commerce-intensive categories, digital services (rebound/direct), with
no-channel categories (rent, insurance) as placebos. Category mapping
justification (see 00\_constants.R).}

\subsection{Mechanism}
\todo{Out-of-home leisure frequency (endline E1) mediating the transport
gradient.}

\subsection{Heterogeneity}
\todo{Pre-specified: age, gender, urbanity.}

\subsection{Design caveats}
\todo{Timing: screen time measured at endline, consumption observed 2019--2024;
reverse-causality/stability argument.}

\section{Results}

\subsection{Screen-time validity}
\todo{Device-assisted screen time vs.\ Swedish benchmarks (Internetstiftelsen,
Svenskarna och internet). See output/screen\_time\_validity.csv.}

\subsection{The net gradient}
\todo{}

\subsection{Decomposition across consumption categories}
\todo{}

\subsection{Mechanism: out-of-home leisure}
\todo{}

\subsection{Heterogeneity}
\todo{}

\section{Discussion}

\todo{Implications for scenario and policy work on low-carbon digitalization;
parameters for prospective life-cycle scenarios; methodological template for
linking financial-transaction and survey data.}

\section*{Pre-registration}
\todo{Link to OSF pre-registration (required before touching real data in the
TRE).}

% \bibliographystyle{plainnat}
% \bibliography{references}

\appendix

\section{Survey instruments}
\todo{Baseline q11/q11b batteries; endline E1, E4 wording.}

\section{Robustness}
\todo{}

\end{document}
